\documentclass[12pt]{article}
\usepackage{amsmath}
\begin{document}

\textit{Proposed Solution to \#5823 SSMJ }\\ \\
\textit{Solution proposed by Prakash Pant, The University of Vermont, Bardiya, Nepal.}\\ \\
\textit{Problem proposed by Daniel Sitaru, National Economic College "Theodor Costescu"
Drobeta Turnu - Severin, Romania
} \\ \\
\textbf{Statement of the Problem:}
Prove that if \( 0<a \le b\), then:
\[ \left(\int_a^b e^{-x^2} dx \right)^2 \ge \left( \int_{\frac{a+b}{2}}^{b}e^{-x^2} dx + \int_{\sqrt{ab}}^{b} e^{-x^2} dx \right)  \left( \int_{a}^{\frac{a+b}{2}}e^{-x^2} dx + \int_{a}^{\sqrt{ab}} e^{-x^2} dx \right)  \]

\textbf{Solution of the Problem:} \\
Let F(x) be the real antiderivative of \( e^{-x^2}\). Then the given problem can be restated as: 
\[ [F(b)-F(a)]^2 \ge \left[  F(b) - F(\frac{a+b}{2}) + F(b) - F(\sqrt{ab}) \right] \left[ F(\frac{a+b}{2})-F(a)+F(\sqrt{ab})-F(a) \right] \]

\[ [F(b)-F(a)]^2 \ge \left[ 2F(b) - F(\frac{a+b}{2}) - F(\sqrt{ab}) \right] \left[ F(\frac{a+b}{2})+F(\sqrt{ab})-2F(a) \right] \]

\[ F^2(b)-2F(b)F(a)+F^2(a) \ge 2F(b)\left[F(\frac{a+b}{2})+F(\sqrt{ab}) \right] - 4 F(b)F(a) -\left[F(\frac{a+b}{2})+F(\sqrt{ab}) \right]^2 \] \[ 
 + 2 F(a) \left[F(\frac{a+b}{2}) + F(\sqrt{ab}) \right] \] 
which can be rewritten as 
\[ \left[F(b)+F(a)\right]^2 + \left[F(\frac{a+b}{2})+F(\sqrt{ab}) \right]^2 \ge 2 \left[F(\frac{a+b}{2}) + F(\sqrt{ab}) \right] \left[F(b) + F(a) \right] 
\]
\[ \left[F(b)+F(a)-F(\frac{a+b}{2})-F(\sqrt{ab}) \right]^2 \ge 0  \]
Since square of a real number cannot be negative, the given statement is true which proves the original statement we began with, hence concluding the proof. 






 
\end{document}